%\newsection{Kapitel1}

%Hier den Inhalt mit \input einfügen und die folgenden Zeilen entfernen

\newvideofile{03KapazitaetundKondensatoren}{Kapazität und Kondensator}
%%%%%%%%%%%%%%%%%%%%%%%%%%%%%%%%%%%%%%%%%%%%%%%%% Anfang - Kapazität	 %%%%%%%%%%%%%%%%%%%%%%%%%%%%%%%%%%%%%%%%%%%%%%

\section{Kapazität und Kondensator}
\s{
	In der Elektrotechnik gibt es drei grundlegende Komponenten, mit denen das elektrische Verhalten von Bauteilen beschrieben werden kann.
	Neben dem zuvor behandelten ohmschen Widerstand, gibt es noch die Kapazität und die Induktivität.
	In diesem Abschnitt geht es um die Eigenschaft der Kapazität, sowie dem dazugehörigen Bauteil, dem Kondensator.
	Der darauffolgende Abschnitt behandelt die Induktivität sowie das zugehörige Bauteil, die Spule.}

	%%%%%%%%%%%%%%%%%%%%%%%%%%%%%%%%%%%%%%%%%%%%%%%%% Anfang - Lernziele Kapazität	 %%%%%%%%%%%%%%%%%%%%%%%%%%%%%%%%%%%%%%%%%%%%%%
	% \begin{Lernziele}{Kapazität und Kondensator}

	% 	Die Studierenden können
	% 	\begin{itemize}
	% 		\item zwischen Kapazität und Kondensator differenzieren.
	% 		\item die wichtigsten Parameter rund um den Kondensator und die Kapazität benennen.
	% 		\item die Kapazität eines Kondensators berechnen.
	% 	\end{itemize}

	% \end{Lernziele}
% \b{
% 	\begin{frame}
% 		\ftx{Lernziele}
% 		\begin{Lernziele}{Kapazität und Kondensator}

% 			Die Studierenden können
% 			\begin{itemize}
% 				\item zwischen Kapazität und Kondensator differenzieren.
% 				\item die wichtigsten Parameter rund um den Kondensator und die Kapazität benennen.
% 				\item die Kapazität eines Kondensators berechnen.
% 			\end{itemize}
% 		\end{Lernziele}
% 	\end{frame}
% }
\v{
	\begin{frame}
		\ftx{Lernziele}
		\begin{Lernziele}{Kapazität und Kondensator}

			Du kannst
			\begin{itemize}
				\item zwischen Kapazität und Kondensator differenzieren.
				\item die wichtigsten Parameter rund um den Kondensator und die Kapazität benennen.
				\item die Kapazität eines Kondensators berechnen.
				\item die Lade- und Entladeeigenschaften erläutern.
			\end{itemize}
		\end{Lernziele}

		\speech{03_Kapazitaet_und_Kondensator1}{1}{% sehr langsame sprache -> neuer seed
			Im nun folgenden Abschnitt beschäftigen wir uns mit dem Begriff der Kapazität und dem damit verbundenen Bauteil dem Kondensator.
			Am Ende dieses Abschnittes solltest du zwischen Kapazität und Kondensator differenzieren können.
			Du kennst die wichtigsten Parameter rund um den Kondensator und die Kapazität, du kannst die Kapazität eines einfachen Plattenkondensators berechnen und abschließend die Lade- und Entladeeigenschaften erläutern.}
	\end{frame}
}
%%%%%%%%%%%%%%%%%%%%%%%%%%%%%%%%%%%%%%%%%%%%%%%%% Ende - Lernziele Kapazität	 %%%%%%%%%%%%%%%%%%%%%%%%%%%%%%%%%%%%%%%%%%%%%%

\subsection{Die elektrische Kapazität $C$}
\s{
	Die elektrische Kapazität ist die Fähigkeit eines Objektes, elektrische Energie in Form eines elektrischen Feldes zu speichern.
	Diese Eigenschaft kommt zustande, sobald elektrische Ladungsträger ungleich verteilt sind und sich zwischen den Bereichen
	ein elektrisches Feld ausbildet. Das Formelzeichen der Kapazität ist $C$ (capacity) während die Einheit in Farad (F) angegeben wird.
    Das einfachste und geläufigste Beispiel, an dem die Kapazität erklärt werden kann, ist der Plattenkondensator.
	Bei dieser Art des Kondensators liegen sich zwei leitfähige Platten gegenüber, die jeweils mit einem anderen Potential beaufschlagt werden.
	In Abbildung \ref{fig:idealer Kondensator} ist dieser mit einem idealisierten Feldverlauf dargestellt.
}
%%%%%%%%%%%%%%%%%%%%%%%%%%%%%%%%%%%%%%%%%%%%%%%%% Anfang - Grafik idealer Kondensator	 %%%%%%%%%%%%%%%%%%%%%%%%%%%%%%%%%%%%%%%%%%%%%%
\begin{frame}[t]
	\ftx{Der Kondensator}
	\b{
		\begin{itemize}
			\item Der ideale Plattenkondensator
		\end{itemize}
	}
	\begin{figure}[H]
		\centering
		\includesvg[width=0.6\textwidth]{Idealer_Plattenkondensator}
		\s{\caption{\textbf{Plattenkondensator mit idealisiertem Feldverlauf.} Der Kondensator speichert elektrische Energie im elektrischen Feld, das zwischen den Platten idealisiert als homogen dargestellt ist.}}
		% und hat ein Kapazität \( C \),von der Plattenfläche, dem Plattenabstand und dem Dielektrikum zwischen den Platten abhängt.
		\label{fig:idealer Kondensator}
		\alttext{Schematische Darstellung eines Plattenkondensators. Links ist eine Platte mit positiven Ladungen markiert, rechts eine Platte mit negativen Ladungen. Zwischen den Platten zeigen mehrere parallele rote Pfeile von der positiven zur negativen Platte und stellen ein homogenes elektrisches Feld dar.}
	\end{figure}
	%%%%%%%%%%%%%%%%%%%%%%%%%%%%%%%%%%%%%%%%%%%%%%%%% Ende - Grafik idealer Kondensator	 %%%%%%%%%%%%%%%%%%%%%%%%%%%%%%%%%%%%%%%%%%%%%%
	\s{
		Im idealen Fall wird angenommen, dass sich außerhalb der Kondensatorplatten kein Feld ausbildet.
		In der Realität sammeln sich an den seitlichen Rändern der Platten ebenfalls Ladungsträger, welche zu komplizierteren Feldverteilungen führen.
		In Abbildung \ref{fig:realer Kondensator} ist die Grafik eines solchen realen Plattenkondensators dargestellt.
	}% nur im Skript

	\speech{03_Kapazitaet_und_Kondensator2}{1}{
		Die einfachste Form eines Kondensators ist der sogenannte Plattenkondensator.
		Bei dem hier dargestellten idealen Plattenkondensator stehen sich zwei leitfähige Platten parallel gegenüber.
		Zwischen diesen leitfähigen Platten ist ein Isolator ein sogenanntes Dielektrikum.
		Das heißt, es kann kein Strom zwischen diesen beiden Platten fließen.
		Wenn es uns nun zum Beispiel durch das Anlegen einer Spannung an die beiden Platten gelingt, eine Ladungsträgertrennung zu erzeugen
		hier dargestellt durch positive Ladungen auf der linken Seite und negative Ladungen auf der rechten Seite
		dann baut sich ein elektrisches Feld E auf, das in der Darstellung in rot eingezeichnet ist.
	}

\end{frame}
%%%%%%%%%%%%%%%%%%%%%%%%%%%%%%%%%%%%%%%%%%%%%%%%% Anfang - Grafik realer Kondensator	 %%%%%%%%%%%%%%%%%%%%%%%%%%%%%%%%%%%%%%%%%%%%%%
\begin{frame}[t]
	\ftx{Der Kondensator}
	\b{
		\begin{itemize}
			\item Der reale Plattenkondensator
		\end{itemize}
	}
	\begin{figure}[H]
		\centering
		\includesvg[width=0.6\textwidth]{Realer_Plattenkondensator}
		\s{\caption{\textbf{Plattenkondensator mit realem Feldverlauf.} Im Gegensatz zum idealen Kondensator treten Randfelder und Verluste auf, die die tatsächliche Kapazität und das elektrische Verhalten beeinflussen.}}
		\alttext{Schematische Darstellung eines Plattenkondensators mit realem Feldverlauf. Links ist eine positiv geladene Platte, rechts eine negativ geladene Platte dargestellt. Zwischen den Platten verlaufen mehrere gerade Feldlinien von der positiven zur negativen Platte. An den oberen und unteren Plattenrändern biegen die Feldlinien nach außen ab und bilden Randfelder, die als gekrümmte Pfeile außerhalb des Plattenbereichs dargestellt sind.}
		\label{fig:realer Kondensator}
	\end{figure}

	\speech{03_Kapazitaet_und_Kondensator3}{1}{
		In der Realität ist das elektrische Feld jedoch nicht nur exakt zwischen den Platten vorhanden und senkrecht.
		Es gibt darüber hinaus sogenannte Streufelder, die an den Rändern der Platten auftreten.
		Da die Platten aus leitfähigem Material bestehen, müssen sämtliche Feldlinien die austreten bzw. eintreten, senkrecht auf den entsprechenden Oberflächen stehen.
		Deshalb entstehen gebogene Feldlinienverläufe am Rand.
		Auch die oben und unten angedeuteten Feldlinien sind geschlossen
		sie sind hier zur besseren Lesbarkeit nicht vollständig dargestellt, weil das sonst die Übersichtlichkeit zerstören würde.
	}

\end{frame}
%%%%%%%%%%%%%%%%%%%%%%%%%%%%%%%%%%%%%%%%%%%%%%%%% Ende - Grafik realer Kondensator	 %%%%%%%%%%%%%%%%%%%%%%%%%%%%%%%%%%%%%%%%%%%%%%
\begin{frame}
	\ftx{Die elektrische Kapazität}
	\b{
		\begin{itemize}\centering
			\item $ \displaystyle C = \frac{Q}{U} \qquad [C] =  1\,\text{Farad} = 1\, \text{F} = 1\,\frac{\text{C}}{\text{V}} = 1\,\frac{\text{A} \text{s}}{\text{V}} \label{eq:cqu_2} $
		\end{itemize}
	}
	\speech{03_Kapazitaet_und_Kondensator4}{1}{
		Ein solcher Plattenkondensator baut ein elektrisches Feld auf.
		In diesem Feld kann elektrische Energie gespeichert werden.
		Das Maß für die Fähigkeit, elektrische Energie im Feld eines Kondensators zu speichern, ist die Kapazität.
		Die Kapazität wird mit dem Buchstaben C abgekürzt.
		Die Einheit der Kapazität ist das Farad.
		Ein Farad bedeutet, dass ein Coulomb an Ladung gespeichert wird, wenn 1 Volt Spannung anliegt.
		Formelmäßig also: C ist gleich Q durch U.
		Alternativ in S-I-Einheiten: Farad ist gleich Ampere-sekunde pro Volt.
	}
	\silence{4}
\end{frame}
\s{
	In vielen Fällen kann das reale Verhalten jedoch vernachlässigt werden, was die Berechnungen deutlich vereinfacht.
	Aus dem Grund wird im Folgenden nur noch von idealisierten Kondensatoren ausgegangen.

	Die Kapazität ist definiert als das Verhältnis der gespeicherten elektrischen Ladung $Q$ zur angelegten Spannung $U$.
	Je höher die Kapazität eines Kondensators ist, desto mehr Ladung kann er bei einer gegebenen Spannung speichern.
	
	\begin{equation}
		C = \frac{Q}{U} \qquad [C] =  1\,\text{Farad} = 1\,\text{F} =  1\,\frac{\text{C}}{\text{V}} = 1\,\frac{\text{A} \text{s}}{\text{V}} \label{eq:cqu_2}
	\end{equation}
	% %%% aus Einführung in zeitabhängige Vorgänge (MH)
	% \begin{equation}
	% 	Q = C \cdot U \quad \Rightarrow \quad C = \frac{Q}{U} = \frac{\varepsilon \cdot \iint \vec{E} \cdot \mathrm{d}\vec{A}}{\int \vec{E}\mathrm{d}\vec{s}} \label{eq:qcu2}
	% \end{equation}
}
%%%%%%%%%%%%%%%%%%%%%%%%%%%%%%%%%%%%%%%%%%%%%%%%% Ende - Formeln Kapazität	 %%%%%%%%%%%%%%%%%%%%%%%%%%%%%%%%%%%%%%%%%%%%%%
\subsection{Der Kondensator als Bauelement}

\s{
	Die Eigenschaften der Kapazität werden in verschiedenen Bereichen der Elektrotechnik bewusst genutzt.
	Die Nutzbarmachung dieser Eigenschaften erfolgt mit Hilfe von Kondensatoren. Es gibt sie in verschiedenen Ausführungen.
	Abbildung \ref{fig:Verschiedene Kondensatoren} zeigt einige Beispiele von THT- bis hin zu SMD-Bauteilen.
}

%%%%%%%%%%%%%%%%%%%%%%%%%%%%%%%%%%%%%%%%%%%%%%%%% Anfang - Bild versch. Kondensatoren	%%%%%%%%%%%%%%%%%%%%%%%%%%%%%%%%%%%%%%%%%%%%%%
\begin{frame}[t]
	\ftx{Der Kondensator als Bauelement}
	\b{
		\begin{itemize}
			\item Verschiedene Arten und Bauformen von Kondensatoren
		\end{itemize}
	}
	\begin{figure}[H]
		\centering
		\includegraphics[width=0.7\textwidth]{ko}
		\s{\caption{\textbf{Verschiedene Arten und Bauformen von Kondensatoren.} Darunter Keramik- und Elektrolytkondensatoren. Diese variieren in Größe, Kapazität und Spannung, um unterschiedlichen Anwendungen gerecht zu werden.}}
\alttext{Die Abbildung zeigt ein Foto verschiedener Kondensator-Bauformen. Oben sind mehrere bedrahtete Elektrolytkondensatoren in unterschiedlichen Größen und Farben nebeneinander angeordnet. Rechts oben sind sehr kleine SMD-Kondensatoren zu sehen. Unten links liegt ein axialer Kondensator im Metallgehäuse. In der unteren Reihe sind außerdem mehrere kleine Keramikkondensatoren als Scheiben- oder Tropfenform zu sehen.}
		\label{fig:Verschiedene Kondensatoren}
	\end{figure}

	\speech{03_Kapazitaet_und_Kondensator5}{1}{
		Auf dem hier dargestellten Foto siehst du verschiedene Ausführungen von Bauteilen. Bis auf den ganz rechten Rand sind das alles Durchsteckelemente, also T-H-T-Bauteile, was gerade bei Kondensatoren relativ häufig vorkommt, da diese räumlich groß sein müssen, um hinreichende Kapazitäten zu haben. Wir werden später sehen, dass je größer die Fläche ist, die wir haben, und je kleiner der Abstand, desto größer wird die Kapazität, weil letztlich das elektrische Bild natürlich größer wird und die Ladungsmenge, die gespeichert werden kann, größer wird.
		Weswegen du auch hier sehr schön sehen kannst, dass wir gewisse Zusammenhänge haben. Der rote Kondensator oben links ist ein 400 Volt, 68 mikro Farad Kondensator. Der ist räumlich sehr viel größer als der schwarze 63 Volt, 330 mikro Farad Kondensator rechts daneben. Dieser  hat eine Kapazität, die rund 4,5 bis 5 Mal größer ist als die des roten Kondensators und ist räumlich trotzdem kleiner. Wie kann das erreicht werden? Indem der Abstand zwischen den Platten deutlich kleiner ist, weswegen er auch nur eine kleine Spannung aushalten kann.
		Denn die Platten müssen ja voneinander isoliert sein. Das heißt, man wickelt dort ein in Elektrolyt getränktes Material zwischen die Platten ein. Das sind sogenannte Elektrolytkondensatoren. Das hat nur eine gewisse Dicke und hinzufolge auch nur eine gewisse Spannungsfestigkeit.
		Das heißt, wenn du Kondensatoren haben möchtest, die hohe Spannung aushalten und gleichzeitig aber hohe Kapazitätswerte haben, dann werden das räumlich extrem große Bauteile. Unten ist eine andere Technologie. Das sind Keramik- und Tantalkondensatoren, die meistens deutlich kleinere Kapazitätswerte haben, dafür aber elektrisch häufig sehr viel bessere Eigenschaften haben, sodass die abhängig von der Anwendung gewählt werden. Ganz oben rechts die kleinen Chips. Das sind klassische SMD-Bauteile. Wieder ohne Beinchen, das heißt deutlich weniger parasitäre Effekte, aber auch kleinere Kapazitätswerte. Und unten rechts siehst du Elektrolytkondensatoren als SMD-Kondensatoren, die eben auch direkt auf der Platine montiert werden.
	}
\end{frame}
%%%%%%%%%%%%%%%%%%%%%%%%%%%%%%%%%%%%%%%%%%%%%%%%% Ende - Bild versch. Kondensatoren	 %%%%%%%%%%%%%%%%%%%%%%%%%%%%%%%%%%%%%%%%%%%%%%
\s{
	Darüber hinaus gibt es sie aber auch in deutlich kleineren bzw. deutlich größeren Dimensionen.
	In der Hochfrequenztechnik werden Kapazitäten beispielsweise lediglich durch das Design der Leiterplatten realisiert,
	während Kondensatoren für industrielle oder energietechnische Anwendungen mehrere Zentimeter bis Meter groß werden können.

	Was alle Kondensatoren verbindet, ist ihre Fähigkeit, die Eigenschaft der Kapazität nutzbar zu machen,
	wie es im folgenden Schaltbild \ref{fig:idealer Kondensator_ESB_1} dargestellt ist. Dieses Schaltbild repräsentiert jedoch nur die idealisierte, nutzbare Kapazität.
}
%%%%%%%%%%%%%%%%%%%%%%%%%%%%%%%%%%%%%%%%%%%%%%%%% Anfang - Schaltzeichen Kapazität U,I	%%%%%%%%%%%%%%%%%%%%%%%%%%%%%%%%%%%%%%%%%%%%%%
\begin{frame}[t]
	\ftx{Der Kondensator als Bauelement}
	% idealer Kondensator
	\b{
		\begin{itemize}
			\item Idealer Kondensator
		\end{itemize}
		\begin{figure}[h]
			\begin{center}
				\input{Tikz/src/03_Kapazitaet_Kondensator/Ersatzschaltbild_idealer_Kondensator}
	\alttext{Schaltsymbol eines idealen Kondensators C sind zwei linien senkrecht zur Leitungslinie. Ein Pfeil kennzeichnet die Kondensatorspannung uC von t, und ein Pfeil kennzeichnet den Strom iC von t durch den Kondensator.}
			\end{center}
		\end{figure}

		\only<2->{

			\begin{itemize}
				\item Ersatzschaltbild eines realen Kondensators
			\end{itemize}
			\begin{figure}[H]
				\begin{center}
					\input{Tikz/src/03_Kapazitaet_Kondensator/Ersatzschaltbild_realer_Kondensator}
				\alttext{Ersatzschaltbild eines realen Kondensators als Reihenschaltung. Von links nach rechts sind ein idealer Kondensator C, ein Serienwiderstand Rs und eine Serieninduktivität Ls dargestellt. Der Strom iC von t fließt nach rechts. Die Spannungsabfälle sind als Pfeile markiert: uC von t am Kondensator, uR von t am Widerstand und uL von t an der Induktivität.}	
				\end{center}
			\end{figure}
		}
	}

	\speech{03_Kapazitaet_und_Kondensator6}{1}{
		In einer idealen Welt hat ein Kondensator tatsächlich nur die gewünschte Kapazität. An dieser gewünschten Kapazität liegt eine Spannung an, die jetzt hier als zeitlich variable Spannung anliegt. Und dann fließt auch tatsächlich ein zeitlich variabler Strom durch eben genau diese Kapazität. Ein Strom IC von T.
	}

	\speech{03_Kapazitaet_und_Kondensator6}{2}{
		Leider Gottes schlagen aber die parasitären Elemente bei einem realen Kondensator sehr stark ins Gewicht. Das einfachste Ersatzschaltbild eines realen Kondensators ist die gewünschte Kapazität in Reihe zu einem Serienwiderstand und einer Serieninduktivität. Der Serienwiderstand, den können wir uns relativ leicht erklären, denn auf dem Bild haben wir ja schon gesehen das wir irgendwo immer Anschlussleitungen haben müssen. Diese Anschlussleitungen haben immer einen ohmschen Widerstand der verändert natürlich das Verhalten des Kondensators. Das wäre meistens noch sehr gut zu verschmerzen.
		Viel problematischer ist aber die in Serie liegende Induktivität. Im späteren Verlauf, in anderen Modulen, wirst du das Frequenzverhalten von Kapazität und Induktivität näher kennenlernen. Vereinfacht gesprochen leitet ein Kondensator den Strom immer besser, je höher die Frequenz ist. Wohingegen die Induktivität den Strom immer schlechter leitet mit höher werdender Frequenz.
		Sodass wir hier im Prinzip gegensätzliche Bauteileigenschaften haben. In der Realität ist es tatsächlich so, dass ein Kondensator abhängig von seinem Aufbau irgendwann nicht mehr die gewünschte Kapazität schaltungstechnisch darstellt, sondern sich eher wie eine Induktivität verhält. Da die serielle parasitäre Induktivität dominant wird. Und je größer das Bauelement geometrisch ist, desto schlechter ist dieses Verhältnis zwischen Nennkapazität und parasitärer Induktivität.
		Das heißt die gerade dargestellten Elektrolytkondensatoren, die räumlich sehr groß sind, die kannst du tatsächlich nur bei relativ niedrigen Frequenzen verwenden. Da ist dann der kapazitive Anteil noch dominant. Daher werden diese Kondensatoren z.B. benutzt, um Eingangsspannung zu stabilisieren. Wohingegen für hochfrequentere Anwendungen, Filterschaltung und so, kommen eher die ganz kleinen SMD-Bauteile zum Tragen, weil dann das kapazitive Verhalten des Kondensators bis hin zu höheren Frequenzen garantiert werden kann.
	}

\end{frame}
\s{
	\begin{figure}[H]
		\begin{center}
			\input{Tikz/src/03_Kapazitaet_Kondensator/Ersatzschaltbild_idealer_Kondensator}
			\caption{\textbf{Schaltelement eines idealen Kondensators.} Mit einem Spannungsabfall und einem Stromfluss.}
\alttext{Schaltsymbol eines idealen Kondensators C sind zwei linien senkrecht zur Leitungslinie. Ein Pfeil kennzeichnet die Kondensatorspannung uC von t, und ein Pfeil kennzeichnet den Strom iC von t durch den Kondensator.}
			\label{fig:idealer Kondensator_ESB_1}
		\end{center}
	\end{figure}
	%%%%%%%%%%%%%%%%%%%%%%%%%%%%%%%%%%%%%%%%%%%%%%%%% Ende - Schaltzeichen Kapazität U,I	 %%%%%%%%%%%%%%%%%%%%%%%%%%%%%%%%%%%%%%%%%%%%%%
	Für eine realitätsgetreue Nachbildung der Schaltung ist es wichtig den Kondensator als Bauteil vollständig zu beschreiben.
	Dies wird mit Hilfe eines Ersatzschaltbildes gemacht, welches auch die ungewollten, also parasitären Eigenschaften
	des Bauteils, nachbildet. Abbildung \ref{fig:realer Kondensator_ESB_2} zeigt ein solches Ersatzschaltbild.

	%%%%%%%%%%%%%%%%%%%%%%%%%%%%%%%%%%%%%%%%%%%%%%%%% Anfang - ESB Kondensator	 %%%%%%%%%%%%%%%%%%%%%%%%%%%%%%%%%%%%%%%%%%%%%%
	% realer Kondensator


	\begin{figure}[H]
		\begin{center}
			\input{Tikz/src/03_Kapazitaet_Kondensator/Ersatzschaltbild_realer_Kondensator}
			\caption{\textbf{Ersatzschaltbild eines realen Kondensators.} Der ideale Kondensator $C$ wird durch eine Serien-Induktivität
				$L_\mathrm{s}$ und einen Serien-Widerstand $R_\mathrm{s}$ ergänzt, um parasitäre Effekte zu berücksichtigen.}
\alttext{Ersatzschaltbild eines realen Kondensators als Reihenschaltung. Von links nach rechts sind ein idealer Kondensator C, ein Serienwiderstand Rs und eine Serieninduktivität Ls dargestellt. Der Strom iC von t fließt nach rechts. Die Spannungsabfälle sind als Pfeile markiert: uC von t am Kondensator, uR von t am Widerstand und uL von t an der Induktivität.}
			\label{fig:realer Kondensator_ESB_2}
		\end{center}
	\end{figure}
}

%%%%%%%%%%%%%%%%%%%%%%%%%%%%%%%%%%%%%%%%%%%%%%%%% Ende - ESB Kondensator	 %%%%%%%%%%%%%%%%%%%%%%%%%%%%%%%%%%%%%%%%%%%%%%
\s{
	Im Ersatzschaltbild des Kondensators werden nicht nur die nutzbare Kapazität, sondern auch die parasitären Eigenschaften
	berücksichtigt, die in realen Kondensatoren auftreten. Diese parasitären Eigenschaften entstehen beispielsweise durch die
	Induktivität der Anschlüsse und die ohmschen Verluste im Material.
	Das modellierte Ersatzschaltbild ermöglicht es, diese Effekte im Schaltungsentwurf zu berücksichtigen
	und das Schaltverhalten besser vorherzusagen.
}
%%%%%%%%%%%%%%%%%%%%%%%%%%%%%%%%%%%%%%%%%%%%%%%%%%%% Anfang Merksatz - Kondensator %%%%%%%%%%%%%%%%%%%%%%%%%%%%%%%%%%%%%%%%%
\begin{frame}
	\ftx{Merksatz}
	\begin{Merksatz}{Kondensator}

	Der Kondensator ist der verzweifelte Versuch eine Kapazität nachzubilden.
	\end{Merksatz}

	\speech{03_Kapazitaet_und_Kondensator7}{1}{
		Wir können uns also merken, dass der Kondensator immer nur der verzweifelte Versuch ist, eine Kapazität nachzubilden. Wir hatten das ja schon beim Widerstand, dass das Bauteil der Widerstand eben noch andere Eigenschaften hat, als nur die gewünschte Eigenschaft des  Widerstandes. So ist das letztlich beim Kondensator auch. Und diese ungewünschten Eigenschaften fallen insbesondere dann ins Gewicht, wenn wir zu höheren Frequenzen hingehen.
	}

\end{frame}
%%%%%%%%%%%%%%%%%%%%%%%%%%%%%%%%%%%%%%%%%%%%%%%%%%%% Ende Merksatz - Kondensator %%%%%%%%%%%%%%%%%%%%%%%%%%%%%%%%%%%%%%%%%
%%%%%%%%%%%%%%%%%%%%%%%%%%%%%%%%%%%%%%%%%%%%%%%%% Anfang - Dielektrikum	 %%%%%%%%%%%%%%%%%%%%%%%%%%%%%%%%%%%%%%%%%%%%%%
\subsection{Das Dielektrikum}
\s{
	Das Dielektrium ist, wie der Name bereits impliziert, ein dielektrisches, also elektrisch nicht leitfähiges Material.
	Da jedes Material aus Atomen besteht und alle Atomrümpfe mit Elektronen versehen sind (Elektronenhülle),
	wirkt sich entsprechend auch ein elektrisch nicht leitfähiges Material auf das Verhalten des elektrischen Feldes aus.
	Sobald ein solches Material zwischen die Platten eines Kondensators eingefügt wird, ändern sich demnach auch die Eigenschaften des Kondensators.
	In Abbildung \ref{fig:idealer Kondensator mit Dielektrikum} ist anhand der unterbrochenen Pfeilung des elektrischen Feldes
	das unterschiedliche Verhalten zu erkennen. Die elektrische Flussdichte $D$ ist eine Hilfsgröße, welche materialunabhängig ist,
	auf welche in Abschnitt \ref{subsec:D} noch eingegangen wird.
}
%%%%%%%%%%%%%%%%%%%%%%%%%%%%%%%%%%%%%%%%%%%%%%%%% Anfang - Grafik Dielektrikum	 %%%%%%%%%%%%%%%%%%%%%%%%%%%%%%%%%%%%%%%%%%%%%%
\begin{frame}[t]
	\ftx{Das Dielektrikum}
	\begin{figure}[H]
		\centering
		\includesvg[width=0.6\textwidth]{Idealer_Plattenkondensator_Dielektrikum}
		\s{\caption{\textbf{Idealer Kondensator mit Dielektrikum.} Der Kondensator besteht aus zwei Platten, zwischen denen sich ein Dielektrikum befindet. Das Dielektrikum
				ist ein nicht leitendes Material, das die elektrische Feldstärke reduziert und die Kapazität des Kondensators erhöht, indem es die Relative Dielektrizitätskonstante erhört.}}
		\alttext{Schematische Darstellung eines Plattenkondensators mit Dielektrikum. Links ist die positiv geladene Platte mit Pluszeichen dargestellt, rechts die negativ geladene Platte mit Minuszeichen. Zwischen den Platten befindet sich ein graues, mittig eingezeichnetes Dielektrikum. Mehrere horizontale Pfeile zeigen die Feldlinien von der positiven zur negativen Platte. In der Mitte sind das elektrische Feld E und die elektrische Flussdichte D eingezeichnet. Über der Anordnung steht die Beschriftung Dielektrikum.}
		\label{fig:idealer Kondensator mit Dielektrikum}
	\end{figure}
	\b{\begin{itemize}
  \item Nicht leitendes Material zwischen den Kondensatorplatten
  \item Wird im elektrischen Feld polarisiert
  \item Erhöht die Kapazität
  \item Reduziert die elektrische Feldstärke bei gleicher Ladung
\end{itemize}}
	%;Volatisationseffekt -> Polarisationseffekte
	\speech{03_Kapazitaet_und_Kondensator8}{1}{
		Ich hatte ja bereits gesagt, dass die beiden Platten voneinander isoliert werden müssen. Und dies kann im einfachsten Fall mit Luft als Dielektrikum passieren. Es kann aber auch ein dielektrischer Stoff eingeführt werden, das sogenannte Dielektrikum, hier in dunkelgrau dargestellt. Wir haben dann durchgehend gleich, sowohl im Dielektrikum als auch in der Luft, die elektrische Flussdichte D. Daraus ergibt sich als nicht-kontinuierlichen Teil die elektrische Feldstärke E, die im Dielektrikum deutlich größer werden wird als außerhalb. Was erreichen wir damit? Wir erreichen damit letztlich, dass die Kapazität unseres Kondensators sich erhöht, nämlich durch Polarisationseffekte und ähnliches.
	}

\end{frame}
% Vielleicht an dieser Stelle Einschub zu Dielektrikum,
% Permeativität, was das ist, Zusammenhang zum Wellenwiderstand und zur Lichtgeschwindigkeit.

%%%%%%%%%%%%%%%%%%%%%%%%%%%%%%%%%%%%%%%%%%%%%%%%% Ende - Grafik Dielektrikum	 %%%%%%%%%%%%%%%%%%%%%%%%%%%%%%%%%%%%%%%%%%%%%%
\s{
	Die Eigenschaften des Dielektrikums sind materialspezifisch und werden durch die Dielektrizitätszahl $\varepsilon_\mathrm{r}$ quantifiziert.
	In Tabelle \ref{tab:dielektrizität} sind Beispiele einiger Dielektrika und deren $\varepsilon_\mathrm{r}$, aufgeführt:
}
%%%%%%%%%%%%%%%%%%%%%%%%%%%%%%%%%%%%%%%%%%%%%%%%% Anfang - Tabelle epsylon	 %%%%%%%%%%%%%%%%%%%%%%%%%%%%%%%%%%%%%%%%%%%%%%
\begin{frame}[t]
	\ftx{Das Dielektrikum}
	\b{
		\begin{itemize}
			\item Relative Dielektrizitätszahl verschiedener Materialien
		\end{itemize}
		\vspace{1cm}
	}
	%%%%% Dielektrizitätszahl Einfache Tabelle
	\begin{table}[H]
		\centering
		\resizebox{0.3\textwidth}{!}{%
			\begin{tabular}{lc}
				\toprule
				\textbf{Material} & \textbf{\boldmath$\varepsilon_\mathrm{r}$} \\
				\midrule
				Vakuum            & 1                                          \\
				Luft (bei STP)    & 1.0006                                     \\
				Kunststoff (PE)   & 2.25-2.3                                   \\
				Glas              & 3-10                                       \\
				Keramik           & 5-300                                      \\
				Silizium          & 11.7                                       \\
				Tantaloxid        & 25-30                                      \\
				Wasser            & 81                                         \\
				\bottomrule
			\end{tabular}%
		}
		\s{\caption{\textbf{Relative Dielektrizitätszahl verschiedener Materialien.} Werte von $\varepsilon_r$ im Vergleich zum Vakuum.}}
		\alttext{Tabelle mit der relativen Dielektrizitätszahl eps r für verschiedene Materialien. Aufgelistet sind: Vakuum 1, Luft bei STP 1,0006, Kunststoff PE 2,25 bis 2,3, Glas hat 3 bis 10, Keramik 5 bis 300, Silizium 11,7, Tantaloxid 25 bis 30 und Wasser 81.}
		\label{tab:dielektrizität}
	\end{table}

	\speech{03_Kapazitaet_und_Kondensator9}{1}{
		Wir geben diese Möglichkeit der Erhöhung des elektrischen Feldes in einer elektrischen Flussdichte über die Dielektrizitätszahlen an. Die relative Dielektrizitätszahl gibt quasi an, wie viel das elektrische Feld verstärkt wird. Vakuum hat eine Dielektrizitätszahl von eins und Luft hat ebenfalls eine von etwa eins. Und der häufig verwendete Kunststoff PE liegt irgendwo bei zwei komma 3. Aber wir können auch Keramiken finden, die eine Dielektrizitätszahl von bis zu 300 haben oder zum Beispiel Tantal-oxid welches ein epsylon-er von fünfundzwanzig bis dreizig hat. Wasser ist ebenfalls ein sehr gutes Dielektrikum. Das Problem bei Wasser ist nur, dass es bei kleinsten Verunreinigungen bereits elektrisch leitfähig wird und damit diese dielektrischen Eigenschaften verliert.
	}
	\silence{4}
\end{frame}

%%%%% Kästchen Tabelle

% \begin{table}[H]
%     \centering
%     \begin{tabular}{|l|c|}
%         \hline
%         \textbf{Material} & \textbf{\boldmath$\varepsilon_\mathrm{r}$} \\
%         \hline
%         Vakuum & 1 \\
%         \hline
%         Luft (bei STP) & 1.0006 \\
%         \hline
%         Kunststoff (PE) & 2.25-2.3 \\
%         \hline
%         Glas & 3-10 \\
%         \hline
%         Keramik & 5-300 \\
%         \hline
%         Silizium & 11.7 \\
%         \hline
%         Tantaloxid & 25-30 \\
%         \hline
%         Wasser & 81 \\
%         \hline
%     \end{tabular}
%     \caption{Dielektrizitätszahl verschiedener Materialien}
%     \label{tab:dielektrizität}
% \end{table}
%%%%%%%%%%%%%%%%%%%%%%%%%%%%%%%%%%%%%%%%%%%%%%%%% Ende - Tabelle epsylon	 %%%%%%%%%%%%%%%%%%%%%%%%%%%%%%%%%%%%%%%%%%%%%%
\s{
	Zur vollständigen Beschreibung der Materialeigenschaften des Dielektrikums wird neben der Dielektrizitätszahl noch die elektrische Feldkonstante $\varepsilon_\mathrm{0}$ benötigt.
	Dabei handelt es sich um das dielektrische Verhalten des elektrischen Feldes im Vakuum. Während die Dielektrizitätszahl $\varepsilon_\mathrm{r}$ einheitslos ist,
	bringt die elektrische Feldkonstante die Einheit $\frac{\text{As}}{\text{Vm}}$ mit. Sie wird mit der relativen Dielektrizitätszahl multipliziert was zusammen
	die Dielektrizitätskonstante $\varepsilon$ ergibt.
}
%%%%%%%%%%%%%%%%%%%%%%%%%%%%%%%%%%%%%%%%%%%%%%%%%%%% Anfang Formel epsylon und Legende %%%%%%%%%%%%%%%%%%%%%%%%%%%%%%%%%
\begin{frame}
	\ftx{Die Dielektrizitätskonstante $\varepsilon$}
	\b{
		\begin{equation*}
			\boldsymbol\varepsilon = \boldsymbol\varepsilon_\mathrm{0} \cdot \boldsymbol\varepsilon_\mathrm{r} \qquad [\boldsymbol\varepsilon] = 1\,\frac{\text{Farad}}{\text{m}} = 1\,\frac{\text{F}}{\text{m}} = 1\,\frac{\text{As}}{\text{Vm}}
		\end{equation*}
		\vspace{0.1cm}
		\begin{flalign*}
	\boldsymbol{\varepsilon}        &:\ \text{Dielektrizitätskonstante } [\text{As}/\text{Vm}] &\\
	\boldsymbol{\varepsilon}_0      &:\ \text{Elektrische Feldkonstante } [8.85421878\cdot 10^{-12}\,\text{As}/\text{Vm}] &\\
	\boldsymbol{\varepsilon}_r      &:\ \text{Relative Dielektrizitätszahl} &\\
	\end{flalign*}
	}

	\speech{03_Kapazitaet_und_Kondensator10}{1}{
		Wofür brauchen wir jetzt die relative Dielektrizitätszahl? Die brauchen wir zum Berechnen von der Dielektrizitätskonstante Epsilon. Diese hat die Einheit Farad pro Meter, denn sie gibt letztlich an, wie stark das Dielektrikum das elektrische Feld verstärkt. Und sie setzt sich zusammen aus Epsilon null und Epsilon r. Epsilon null ist die elektrische Feldkonstante mit einem Wert von acht komma acht fünf fünf mal 10 hoch minus 12 Farad pro Meter oder Ampere-Sekunden pro Volt-Meter. Und Epsilon r ist die einheitenlose elektrische Dielektrizitätszahl.
	}

\end{frame}
\s{
	\begin{equation}
			\boldsymbol\varepsilon = \boldsymbol\varepsilon_\mathrm{0} \cdot \boldsymbol\varepsilon_\mathrm{r} \qquad [\boldsymbol\varepsilon] = 1\,\frac{\text{Farad}}{\text{m}} = 1\,\frac{\text{F}}{\text{m}} = 1\,\frac{\text{As}}{\text{Vm}}
	\end{equation}
	\vspace{0.1cm}
		\begin{flalign*}
	\boldsymbol{\varepsilon}        &:\ \text{Dielektrizitätskonstante } [\text{As}/\text{Vm}] &\\
	\boldsymbol{\varepsilon}_0      &:\ \text{Elektrische Feldkonstante } [8.85421878\cdot 10^{-12}\,\text{As}/\text{Vm}] &\\
	\boldsymbol{\varepsilon}_r      &:\ \text{Relative Dielektrizitätszahl} &\\
	\end{flalign*}
	%%%%%%%%%%%%%%%%%%%%%%%%%%%%%%%%%%%%%%%%%%%%%%%%% Ende Formel epsylon und Legende %%%%%%%%%%%%%%%%%%%%%%%%%%%%%%%%%

	Mit Hilfe dieser Materialeigenschaften des Dielektrikums und den Dimensionen sowie der Entfernung der Kondensatorplatten,
	kann die Kapazität des Kondensators wie folgt berechnet werden:
}
\begin{frame}
	\ftx{Die Kapazität $C$}
	\b{
		\begin{equation*}
				C = \frac{\boldsymbol\varepsilon \cdot A}{d} \qquad [C] = 1\,\text{F}
		\end{equation*}
		\vspace{0.1cm}
		\begin{flalign*}
			A &:\ \text{Querschnittsfläche des Materials } [\text{m}{^2}] &\\
			d &:\ \text{Abstand der Elektroden } [\text{m}]               &\\
		\end{flalign*}

	}
	\s{
		%%%%%%%%%%%%%%%%%%%%%%%%%%%%%%%%%%%%%%%%%%%%%%%%%%%% Anfang Formel C (epsilon) und Legende %%%%%%%%%%%%%%%%%%%%%%%%%%%%%%%%%
		\begin{equation}
			C = \frac{\boldsymbol\varepsilon \cdot A}{d} \qquad [C] = 1\,\text{F}
		\end{equation}
		\vspace{0.1cm}
		\begin{flalign*}
			A &:\ \text{Querschnittsfläche des Materials } [\text{m}{^2}] &\\
			d &:\ \text{Abstand der Elektroden } [\text{m}]                &\\
		\end{flalign*}}
	%%%%%%%%%%%%%%%%%%%%%%%%%%%%%%%%%%%%%%%%%%%%%%%%%%%% Ende Formel C (epsilon) und Legende %%%%%%%%%%%%%%%%%%%%%%%%%%%%%%%%%
	%%%%%%%%%%%%%%%%%%%%%%%%%%%%%%%%%%%%%%%%%%%%%%%%% Anfang - el. Flussdichte	 %%%%%%%%%%%%%%%%%%%%%%%%%%%%%%%%%%%%%%%%%%%%%%
	\speech{03_Kapazitaet_und_Kondensator11}{1}{
		Wie kann man jetzt die Kapazität eines Kondensators, die ja in Farad angegeben wird, berechnen? Wir haben ja schon gesehen, dass je größer die Fläche A ist und je kleiner der Plattenabstand d, sich die Kapazität verändert. Genau das ist auch der Zusammenhang: C ist gleich Epsilon mal A durch D. Durch das Einbringen eines Dielektrikums mit hoher Dielektrizitätszahl kann die Kapazität also bei gleichen geometrischen Abmessungen erhöht werden. Und das ist oft erstrebenswert.
	}

\end{frame}


\subsection{Die elektrische Flussdichte}\label{subsec:D}

\begin{frame}[t]
	\ftx{Die elektrische Flussdichte $\vec{D}$}
	\b{
		\begin{equation*}
			\vec{D} = \boldsymbol\varepsilon \cdot \vec{E} + P
		\end{equation*}

		\only<2->{\begin{itemize}
				\item Im statischen Fall kann die Polarisation $P$ oft vernachlässigt werden
			\end{itemize}
			\vspace{2cm}
		}

		\only<3->{
			\begin{equation*} %; done TODO formel Wort Coulomb durch C ersetzen und gespeochenene text anschauen (herleitung von As/m^2 zeigen)
				%[D] = 1\,\frac{\text{Coulomb}}{\text{m}^2} = 1\,\frac{\text{C}}{\text{m}^2} = 1\,\frac{\text{As}}{\text{m}^2}
				[D] = 1\,\frac{\text{C}}{\text{m}^2} = 1\,\frac{\text{As}}{\text{m}^2}
			\end{equation*}
			\vspace{0.1cm}
			\begin{equation*}
				\vec{D} = \boldsymbol\varepsilon \cdot \vec{E}
			\end{equation*}}
	}

	\speech{03_Kapazitaet_und_Kondensator12}{1}{
		Wie kann jetzt die elektrische Flussdichte bestimmt werden? Die elektrische Flussdichte D ist gleich Epsilon mal E plus P. P steht für die Polarisationseffekte, Epsilon ist die Dielektrizitätskonstante, E ist die elektrische Feldstärke.
	}
	\silence{2}
	\speech{03_Kapazitaet_und_Kondensator12}{2}{
		Im statischen Fall, das heißt, wir betrachten erstmal nur Gleichspannung und Gleichstrom, kann die Polarisation P eigentlich vernachlässigt werden.
	}

	\speech{03_Kapazitaet_und_Kondensator12}{3}{
		Die elektrische Flussdichte D hat dann die Einheit Coulomb pro Quadratmeter. Denn wir haben die elektrische Feldstärke in Volt pro Meter und Epsilon, also die Dielektrizitätskonstante, in Farad pro Meter – also Ampere-Sekunden pro Volt-Meter. Multipliziert man das, ergibt sich für D die Einheit Ampere-Sekunden pro Quadratmeter, also Coulomb pro Quadratmeter. Der Zusammenhang lautet also D ist gleich Epsilon mal E.
	}
	\silence{4}
\end{frame}
\s{
	Die elektrische Flussdichte $\vec{D}$, auch Verschiebungsdichte genannt, steht proportional zum elektrischen Feld $\vec{E}$.
	Ihr Vorteil besteht darin, dass sie materialunabhängig ist, was mathematische Vorteile mit sich bringt.

	\begin{equation}
		\vec{D} = \boldsymbol\varepsilon \cdot \vec{E} + P
	\end{equation}

	Zusätzlich zur Dielektrizitätskonstanten $\varepsilon$ wird die Polarisation des Mediums $P$ addiert. Sie beschreibt das Polarisationsverhalten, welches im Einschaltvorgang abläuft.
	Im statischen Fall kann die Polarisation $P$ oft vernachlässigt werden, was die Gleichung wie folgt vereinfacht.

	\begin{equation*}
		[D] = 1\,\frac{\text{Coulomb}}{\text{m}^2} = 1\,\frac{\text{C}}{\text{m}^2} = 1\,\frac{\text{As}}{\text{m}^2}
	\end{equation*}
	\vspace{0.1cm}
	\begin{equation}
		\vec{D} = \boldsymbol\varepsilon \cdot \vec{E}
	\end{equation}

}

%%%%%%%%%%%%%%%%%%%%%%%%%%%%%%%%%%%%%%%%%%%%%%%%% Ende - el. Flussdichte	 %%%%%%%%%%%%%%%%%%%%%%%%%%%%%%%%%%%%%%%%%%%%%%
%%%%%%%%%%%%%%%%%%%%%%%%%%%%%%%%%%%%%%%%%%%%%%%%% Anfang - Schaltverhalten Kondensator	 %%%%%%%%%%%%%%%%%%%%%%%%%%%%%%%%%%%%%%%%%%%%%%

%%%%%%%%%%%%%%%%%%%%%%%%%%%%%%%%%%%%%%%%%%%%%%%%% Anfang - Mathematik Kondensator	 %%%%%%%%%%%%%%%%%%%%%%%%%%%%%%%%%%%%%%%%%%%%%%
%%%%%%%%%%%%%%%%%%%%%%%%%%%%%%%%%%%%%%%%%%%%%%%%% Anfang - Übergang zu den zeitabhängigen Vorgängen %%%%%%%%%%%%%%%%%%%%%%%%%
\s{
\subsection{Einführung in die zeitabhängigen Vorgänge}

	Bisher ist bekannt, dass elektrische Felder in leitfähigen Materialien einen Strom $I$ und eine Spannung $U$ verursachen.
	Der Strom $I$ und die Spannung $U$ beschreiben die Felder jedoch nur integral, sie werden also als zeitlich konstant angesehen.
	Im Gegensatz zum Widerstand $R$ gibt es jedoch Bauteile (Kondensatoren und Spulen), die ein zeitlich veränderliches Verhalten aufweisen.
	Der ohm'sche Widerstand $R$ alleine reicht somit nicht aus, um die Verhältnisse bei Netzwerken mit zeitlich veränderlichen Vorgängen zu beschreiben.
	Zur Beschreibung dieser zeitabhängigen Verhalten werden die zeitabhängigen Größen $u(t)$ und $i(t)$ verwendet.

	Die Herleitung der Zeitabhängigkeiten kann mathematisch folgendermaßen gezeigt werden.
	Aus der Formel \ref{eq:cqu_2} ist bereits die Grundgleichung des Kondensators bekannt:

	\begin{equation}
		Q = C \cdot U   \label{eq:qcu}
	\end{equation}
	%%% NEU DAZU %%%
	Zur Beschreibung des zeitlichen Verhaltens des Kondensators, wird die Änderung der Ladung über die Zeit betrachtet

	\begin{equation}
		\frac{\mathrm{d}Q}{\mathrm{d}t} = C \cdot \frac{\mathrm{d}U}{\mathrm{d}t}
	\end{equation}

	$\frac{\mathrm{d}Q}{\mathrm{d}t}$ beschreibt die Änderungsrate der Ladung, welche nach dem Grundgesetz der Elektrotechnik dem Strom $i(t)$ entspricht, der in den Kondensator hinein- oder herausfließt.

	%%%%%%%%%%%%%%%%
	% In Modul 2 wurde vermittelt, dass der Stromfluss auch als Änderung der Ladung pro Zeit angesehen werden kann:

	\begin{equation}
		\frac{\mathrm{d}Q}{\mathrm{d}t} = i(t)
	\end{equation}

	Wird nun in der Formel \ref{eq:qcu} die konstante Ladung Q durch die zeitlich veränderliche Größe $i(t)$ ersetzt,
	erhält auch die Spannung eine zeitliche Abhängigkeit.

	\begin{equation}
		i_\mathrm{c}(t) = C \cdot \frac{\mathrm{d}u_\mathrm{c}(t)}{\mathrm{d}t} \label{eq:ict}
	\end{equation}

	Durch mathematische Umstellung nach $u_\mathrm{c}(t)$ kann die Spannung in Abhängigkeit des Stromes und der Kapazität dargestellt werden:

	\begin{equation}
		u_\mathrm{c}(t) = \frac{1}{C} \cdot \int i_\mathrm{c}(t) \mathrm{d}t \label{eq:uci}
	\end{equation}

	Aus der Gleichung \ref{eq:ict} folgt unmittelbar, dass sich die Spannung an einer Kapazität nicht schlagartig ändern kann,
	da dafür ein unendlich hoher Strom erforderlich wäre.

}
% Folie Einführung in die zeitabhängige Vorgänge
\begin{frame}[t]
	\ftx{Einführung in die zeitabhängigen Vorgänge}

\b{
		\begin{itemize}\centering
			\item  $ \displaystyle Q = C \cdot U $\vspace{0.5cm}
			%%% NEU DAZU %%%
			\uncover<2->{
				\item $ \displaystyle \frac{\mathrm{d}Q}{\mathrm{d}t} = C \cdot \frac{\mathrm{d}U}{\mathrm{d}t} $\vspace{0.5cm}
			}
			\uncover<3->{
				\item
				 $ \displaystyle	\frac{\mathrm{d}Q}{\mathrm{d}t} = i(t) $\vspace{0.5cm}	
			}
			\uncover<4->{
				\item
					 $ \displaystyle i_\mathrm{c}(t) = C \cdot \frac{\mathrm{d}u_\mathrm{c}(t)}{\mathrm{d}t} $\vspace{0.5cm}
			}
			\uncover<5->{
				\item
				 $ \displaystyle u_\mathrm{c}(t) = \frac{1}{C} \cdot \int i_\mathrm{c}(t) \mathrm{d}t$\vspace{0.5cm}
			}

		\end{itemize}
		
		\only<6->{

			\vspace{-7cm}
			\begin{tikzpicture}{remember picture,overlay}
				\fill[white,opacity=0.7](current page.south west) rectangle(current page.north east);
			\end{tikzpicture}
			\centering
			\vspace{-9cm}
			\begin{Merksatz}{Kondensator}

				\begin{itemize}
					\item Nur eine zeitveränderliche Spannung an einer Kapazität führt zu einem Stromfluss durch eine Kapazität
					\item Die Spannung an einer Kapazität kann sich nicht sprunghaft ändern
				\end{itemize}
			\end{Merksatz}

		}
	}

	%% die dunkte mit absatz in speech sind dazu da, damit das ende des letzten satzes nicht abgehackt wird
	\speech{03_Kapazitaet_und_Kondensator13}{1}{%guter seed hackt am ende ab
		Wir gehen jetzt einmal kurz in die Zeitabhängigen Vorgänge, um das Verhalten einer Kapazität im Stromkreis besser erklären zu können.
		Als erstes fangen wir mit der Basisformel an, die wir bei der Einführung der Kapazität bereits gesehen haben. Q ist gleich C mal U.
	}
	\silence{3}
	\speech{03_Kapazitaet_und_Kondensator13}{2}{%; TODO feintuning nötiug ende abgehackt, und englische aussprache der buchstaben
		Wenn wir jetzt beide Seiten nach der Zeit T ableiten erhalten wir d-Q nach d-T ist gleich C mal d-U nach d-T.
	}
	\silence{3}
	\speech{03_Kapazitaet_und_Kondensator13}{3}{ %;TODO feintuning nötig
		Die Zeitliche ableitung der Ladung mit der Einheit Ampersekunde ergibt Ampere, also den Strom i von t.
		So erhalten wir d-Q nach d-T ist gleich I von T.
	}
	\silence{3}
	\speech{03_Kapazitaet_und_Kondensator13}{4}{%guter seed
		So dass wir schreiben können, der Strom I-C von T, der durch die Kapazität fließt, ist gleich C mal U-C von T nach d-T.
		Was können wir aus dieser Formel sehen bzw. herleiten?
		Zum einen: Nur eine zeitveränderliche Spannung an einer Kapazität führt zu einem Stromfluss durch eine Kapazität.
		Denn wenn wir jetzt eine Gleichspannung hätten, dann wäre U-C nach d-T ja gleich 0.
		Und demzufolge ist der Strom I-C von T gleich 0.
	}

	\silence{3}
	\speech{03_Kapazitaet_und_Kondensator13}{5}{ % seed semi ok
		Und logischerweise kann man dann die Spannung U-C von T ausrechnen als 1 durch C mal das Integral I-C von T-d-t.
		Das ist nichts anderes als das Umstellen der oberen Formel nach U-C von T.
		Ferner können wir hier an der Formel direkt schon sehen, dass die Spannung sich an der Kapazität nicht sprunghaft ändern kann.
		Das heißt, wenn wir eine sprunghafte Änderung der Spannung an einer Kapazität hätten, dann müsste der Strom dafür unendlich groß werden.
		Und das kann er ja nicht.
	}
	\silence{4}
	\speech{03_Kapazitaet_und_Kondensator13}{6}%guter seed
	{
		Insofern merken wir uns:
		Nur eine zeitveränderliche Spannung an einer Kapazität führt zu einem Stromfluss durch eine Kapazität und die Spannung an einer Kapazität kann sich nicht sprunghaft ändern.
	}

\end{frame}

%%%% ggf weg
%Im Abschnitt \ref{sec:Induktivität} wird außerdem eingeführt, dass ein elektrischer Strom ein magnetisches Feld bedingt.
%Im Modul 2 wurde bereits erklärt, dass in beiden Feldern Energie gespeichert und transferiert werden kann.

%%%%%%%%%%%%%%%%%%%%%%%%%%%%%%%%%%%%%%%%%%%%%%%%% Ende - Übergang zu den zeitabhängigen Vorgängen %%%%%%%%%%%%%%%%%%%%%%%%%
\subsection{Schaltverhalten eines Kondensators}
\s{
	Das Schaltverhalten beschreibt das Verhalten eines Bauteils bei Änderung der angelegten Spannung.
	Je nachdem ob eine Spannungsquelle ein- oder ausgeschaltet wird, verhalten sich Strom und Spannung am Bauteil
	entsprechend seiner Charakteristika. Zur Analyse des zeitlichen Verhaltens von Bauteilen ist zuerst ein Schaltbild notwendig,
	welches die zu analysierende Schaltung visualisiert. Abbildung \ref{fig:Schaltbild_CR} zeigt ein Schaltbild einer Reihenschaltung
	einer Kapazität $C$ mit einem Widerstand $R$, welche über einen Schalter mit einer Spannung $U_\mathrm{q}$ versorgt werden können.

}
%%%%%%%%%%%%%%%%%%%%%%%%%%%%%%%%%%%%%%%%%%%%%%%%% Ende - Mathematik Kondensator	 %%%%%%%%%%%%%%%%%%%%%%%%%%%%%%%%%%%%%%%%%%%%%%

%%%%%%%%%%%%%%%%%%%%%%%%%%%%%%%%%%%%%%%%%%%%%%%%% Anfang - Schaltbild mit Schalter	 %%%%%%%%%%%%%%%%%%%%%%%%%%%%%%%%%%%%%%%%%%%%%%
\begin{frame}[t] %laden eines Kondensators
	\ftx{Schaltverhalten eines Kondensators: Aufladen}
	\b{
		\onslide<2->{\begin{minipage}[t]{0.3\textwidth}
				\begin{itemize}
					\item[\textcolor{black}{\textbullet}] $\quad I_\mathrm{max} = \frac{U_\mathrm{q}}{R}$
				\end{itemize}
			\end{minipage}}
		\vspace{-1.5cm}
		% Laden eines Kondensators
		\begin{figure}[H]
			\centering
			%\begin{subfigure}[h]{0.5\textwidth}
			\centering
			\begin{circuitikz}[scale=0.8]
    \centering
    \draw (0,-0.3) to[V] (0,4.2)
    -- (1,4.2) to[opening normal closed switch, l={Schalter wird bei $t_\mathrm{0} = 5\tau$ geöffnet}] (3,4.2)
    -- (4,4.2);
    \draw (4,4.2)   to [C,i,v, name=C, l={$C$}] (4,2.2);
    \draw (4,0.2) -- (4,-0.2);
    \draw (4,-0.3) -- (0,-0.3);
    \draw (4,2.2)   to [R,v, name=R, l={$R$}] (4,-0.3);
    \draw[-latex, thick, color=cyan] (-0.7,2.55)  to node[midway,left, color=cyan] {$U_\mathrm{q}$}
    (-0.7,1.35);
    \iarrmore{C}{$i(t)$};
    \varrmore{C}{$u_\mathrm{C}(t)$};
    \varrmore{R}{$u_\mathrm{R}(t)$};

\end{circuitikz}

\onslide<2->{
    \begin{tikzpicture}
    \begin{axis} [
            xlabel={Zeit},
            ylabel={\textcolor{voltage} {$u_\mathrm{C}(t)$}\only<3->{/\textcolor{current}{$i(t)$}} } ,
            % ylabel={
            % 	\textcolor{cyan}{$U_\mathrm{q}$}}
            % 	\only<3->{/ \textcolor{voltage}{$u_\mathrm{C}(t)$}}
            % 	\only<4->{/ \textcolor{current}{$i(t)$}},
            xmin=0, xmax=12.5,
            ymin=-0.2, ymax=1.2,
            ytick={0, 0.37, 0.63, 1},
            yticklabels={%
                    0,%
                    37\%,%
                    63\%,%
                    100\%
                },
            %yticklabels={0,37\%, 63\%, 100\%},
            xtick={0, 1, 2, 3, 4, 5, 6, 7, 8, 9, 10, 11, 12},
            xticklabels={0, \(\tau\), \(2\tau\), \(3\tau\), \(4\tau\), \(5\tau\), \(6\tau\), \(7\tau\), \(8\tau\), \(9\tau\), \(10\tau\), \(11\tau\), \(12\tau\)},
            width=10cm,
            height=3.5cm,
            axis lines=middle,
            every axis x label/.style={at={(current axis.right of origin)},anchor=west},
            every axis y label/.style={at={(current axis.above origin)},anchor=south},
            font=\small,
            legend style={font=\scriptsize}
        ]
        \addplot [voltage, domain=0:12, samples=100] {ifthenelse(x<0, 0, 1 - exp(-(x)))};
        \only<3->{\addplot [current, domain=0:12, samples=100] {ifthenelse(x<0, 0,     exp(-(x)))};}

        %Schalter markierung in Graph
        \addplot [cyan,thick, dotted] coordinates {(5,1) (5,0)};

        % Markierung der Zeitkonstanten
        \draw[dashed] (axis cs:1,0) -- (axis cs:1,0.63) node[right] {};
        \draw[dashed] (axis cs:0,0.63) -- (axis cs:1,0.63) node[right] {};
        
        \only<3->{\draw[dashed] (axis cs:1,0) -- (axis cs:1,0.37) node[right] {};
            \draw[dashed] (axis cs:0,0.37) -- (axis cs:1,0.37) node[right] {};}
        \node[right] at (5,0.5) {$t_\mathrm{0}=5\tau$};

        % Punkte zur Hervorhebung
        \addplot[
            only marks,
            mark=*,
            mark options={scale=0.7, fill=black}
        ]  coordinates {
                (1, 0.63)

            };
        \only<3->{ \addplot[
                only marks,
                mark=*,
                mark options={scale=0.7, fill=black}
            ]  coordinates {
                    (1, 0.37)

                };}

    \end{axis}
\end{tikzpicture}
}
			%\end{subfigure}

			%%%%%%%%%%%%%%%%%%%%%%%%%%%%%%%%%%%%%%%%%%%%%%%%% Ende - Schaltbild mit Schalter	 %%%%%%%%%%%%%%%%%%%%%%%%%%%%%%%%%%%%%%%%%%%%%%
			%%%%%%%%%%%%%%%%%%%%%%%%%%%%%%%%%%%%%%%%%%%%%%%%% Anfang - Diagramm mit Schalter	 %%%%%%%%%%%%%%%%%%%%%%%%%%%%%%%%%%%%%%%%%%%%%%
			%\only<2->{
			% \begin{subfigure}[b]{1\textwidth}
			% \centering
			% \input{Tikz/src/03_Kapazitaet_Kondensator/Verlauf_Kapazitaet_Aufladen}
			%\end{subfigure}
		\end{figure}
	}
	\iffalse %Originaler tex
		\speech{03_Kapazitaet_und_Kondensator14}{1}{
			a, wenn wir jetzt betrachten, was passiert denn, wenn wir einen Kondensator, einen Widerstand in Reihe schalten und wir öffnen oder schließen einen Schalter, sodass die Spannungsquelle UQ dieses System entweder mit Spannung versorgt oder nicht. Dann fällt als erstes mal auf, wir hätten ja eigentlich eine plötzliche Veränderung der Spannung – eben durch das schalten des Schalters.
			Der Schalter hier ist am Anfang geschlossen und wird nach einer gewissen Zeit, die hier mit 5 Tau bezeichnet ist, geöffnet, sodass normalerweise wir dann ja irgendwann einen abrupten Spannungsabfall hätten. Wir wissen aber, dass die Spannung in der Kapazität sich nicht plötzlich ändern kann, weil ansonsten der Strom unendlich groß werden müsste. Hier ist der Strom aber begrenzt auf einen maximalen Strom von U-Q durch R, weil maximal ja die Spannung an R anliegen kann, und demzufolge haben wir einen maximalen Strom von I max ist gleich U-q durch R
			%; TODO fehlender tex
			%. Wir wissen aber, dass die Spannung in der Kapazität sich nicht plötzlich ändern kann, weil ansonsten der Strom unendlich groß werden müsste. Hier ist der Strom aber begrenzt auf einen maximalen Strom von U-Q durch R, weil maximal ja die Spannung an R anliegen kann, und demzufolge haben wir einen maximalen Strom von I max ist gleich U-q durch R. %Die Platten des Kondensators Laden sich nun weiter auf, wodurch immer mehr Spannung über den Kondensator abfällt. Das ist auch an der Formel zu sehen. Die Spannung des Kondensator ist das Integral des Stromes durch den Kondensator über die Zeit.Was passiert nun mit der Spannung am Widerstand? Sie wird kleiner, wodurch der Strom durch den Widertsand sinkt, bis er gegen 0 geht.Da es sich hierbei um eine Serienschaltung handelt, fließt der gleiche Strom durch den Kondensator. %; TODO fehlender text weiter anpassen
		}

		\speech{03_Kapazitaet_und_Kondensator14}{2}{ % ; TODO: only 2
			Was passiert jetzt also, wenn wir den Schalter zum Zeitpunkt T = 0 Sekunden schließen? Die Kapazität war vorher nicht geladen, demzufolge ist die Spannung in der Kapazität 0. Das heißt, die gesamte Spannung UQ fällt an R ab, was einen Stromfluss bedeutet, der gleich dem maximalen Strom ist. Dadurch lädt sich aber die Kapazität auf. Es verschieben sich dort Ladungen, sodass sich eine Spannung U-C aufbaut.
			Diese Spannung UC plus die Spannung UR müssen immer UQ ergeben. Das sagt uns dieses geschlossene Ringintegral EDS. Sodass letztlich die Spannung UR in dem Maße kleiner wird, wie die Spannung UC größer wird. Was heißt das aber im Umkehrschluss? Der Strom I, der durch beide Elemente fließt, wird kleiner. Demzufolge verlangsamt sich die Aufladung.

		}

	\fi
	%Aufladen eines Kondensator
	\speech{03_Kapazitaet_und_Kondensator14}{1}{% TODO feintuning nötig
		Ja, wenn wir jetzt betrachten, was passiert denn, wenn wir einen Kondensator und einen Widerstand in Reihe schalten und wir öffnen oder schließen einen Schalter, sodass die Spannungsquelle U-Q dieses System entweder mit Spannung versorgt oder nicht.
		Dann fällt als erstes mal auf, wir hätten ja eigentlich eine plötzliche Veränderung der Spannung, eben durch das schließen des Schalters am Anfang und dem ungeladenen Kondesator mit U-C gleich 0 Volt.
	}
	\silence{4}
	\speech{03_Kapazitaet_und_Kondensator14}{2}{%
		Da aber die Spannung am Kondensator wie erwähnt nicht springen kann, und U-C am Anfang 0 Volt beträgt, der Kondensator war also ungladen, fällt die gesamte Spannung am Widerstand ab. Dadurch fließt am Anfang der Strom I-max ist gleich U-q durch R.
		Mit zunehmender Zeit lädt sich der Kondensator auf, heißt mehr Spannung fällt über den Kondensator ab. Die Spannung U-C steigt. .
	}

	\speech{03_Kapazitaet_und_Kondensator14}{3}{
		Dadurch wird die Spannung am Widerstand kleiner, wodurch der Strom sinkt.
		Bei 5 Tau wird der Schalter geöffnet. In diesem Fall verhindert der offene Schalter einen weiteren Stromfluss. Der Kondensator kann sich also weder auf- noch entladen, und hält daher seine Spannung.
		Deshalb können auch vom Netz getrennte Schaltungen gefährlich sein, weil die verbauten Kondensatoren lange ihre Ladung halten können.
		Was passiert jetzt, wenn wir uns das entladen einmal anschauen?
	}
	% ; TODO: Hillgärtner: Wenn wir aufladen, dann sollten wir das auch beim Schalter berücksichtigen. Also ich würde es so machen: Der Schalter ist geöffnet und wird zum Zeitpunkt T = 0 geschlossen. Und dass er danach nochmal geöffnet wird, würde ich hier völlig rauslassen – auch unten in den Bildern. Es geht ja um das Aufladen. Und das andere ist für die nächste Folie. 

	%Entladen eines Kondensator
	% \speech{03_Kapazitaet_und_Kondensator14}{2}{ % ; TODO: only 2
	% Die Spannungsquelle wird nun gegen einen Kurzschluss getauscht, und bei t gleich 0 wird der Schalter geschlossen.
	% Die Spannung U-Q treibt den Strom I-max kurzzeitig durch die Schaltung und entlädt sich dabei, der Strom I wird kleiner und geht gegen 0.
	% Warum entläd sich der Kondensator? der Strom durch die Schaltung wir von der Spannung U-C also der Ladungsdifferenz der Platten gespeist. Durch den Strom wird die ladungsdifferenz ausgeglichen und die Spannung U-C fällt ab, bis alle Ladungen ausgeglichen sind.

	% Die Fragedie sich hierbei stellt ist, wie lange dauert dieser entlade oder auflade Vorgang. Hierfür wird Tau verwendet. 1 Tau beschreibt die Zeit die die Spannung des Kondensator von 0V bis 63 Prozent benötigt. Oder bis der Strom von I-Max auf 37 Prozent abfällt.
	% Beim entladen ist es ähnlich, 1 tau entspricht der Zeit, bis die Kondensatorspannugn auf 37 Prozent ihres Ursprungswert abfällt.

	% Tau kann aber auch aus der Kapazität und dem Serienwiderstand berechent werden. Tau ist gleich C mal R.
	% }

	% ; TODO: Hillgärtner: Das würde ich so gar nicht schreiben. UQ wird ja nicht geschaltet. UQ bleibt ja immer gleich, sondern der Schalter bleibt gleich. Also UQ würde ich hier gar nicht reinzeichnen. Ich würde nur UC, eventuell noch UR und I von T reinmachen. Wie gesagt, mit dem Hinweis: Der Schalter sei zum Zeitpunkt T = 0 noch geöffnet und wird zum Zeitpunkt T = 0 geschlossen. Allerdings fehlt noch die Voraussetzung, dass UC(0) = 0 V ist. 


	%originaler speech text ersetzt durch text oberhalb von JW



	%\speech{03_Kapazitaet_und_Kondensator14}{3}{ % ; TODO: only 3
	%Wir sehen also, dass sich die Aufladung des Kondensators verlangsamt. Warum?    
	%}

	% \speech{03_Kapazitaet_und_Kondensator14}{4}{ % ; TODO: only 4
	% Weil eben der Strom kleiner wird. Und der Strom wird eben deshalb kleiner, weil der Strom durch die Spannung UR(t) bedingt ist. Und sobald sich UC(t) aufbaut, wird UR(t) kleiner – also wird I(t) kleiner.
	% Sowohl der Strom als auch die Spannung an der Kapazität folgen jetzt hier einer E-Funktion. Und es gibt eben diese Zeitkonstante Tau, die abhängig ist von R und C. Und diese Zeitkonstante Tau gibt uns eben genau an, dass der Kondensator – wenn er vollständig entladen war – auf 63 Prozent aufgeladen ist, und gleichzeitig der aufladende Strom nur noch 37 Prozent seines Maximalwertes hat. 
	%}

\end{frame}


%entladen des kondensator
\begin{frame}
	\ftx{Schaltverhalten eines Kondensators: Entladen}
	\b{
		\vspace{-0.2cm}
		\begin{minipage}[t]{0.3\textwidth}
			\begin{itemize}
				\item[\textcolor{black}{\textbullet}] $ \quad I_\mathrm{max} = \frac{U_\mathrm{q}}{R}$
				\item[\textcolor{black}{\textbullet}] $ \quad \tau = R \cdot C$
			\end{itemize}
		\end{minipage}
		\vspace{-2.0cm}
		\begin{figure}[H]
			\centering
			\input{Tikz/src/03_Kapazitaet_Kondensator/Laden_Kondensator_ESB_Schalter_schliessen.tex}
		\end{figure}
	}

	\speech{03_Kapazitaet_und_Kondensator15}{1}{ % ; TODO: only 2
		Die Spannungsquelle wird nun gegen einen Kurzschluss getauscht, und bei t gleich 0 wird der Schalter geschlossen.
		Die Kondensatorspannung U-C treibt den Strom I-max kurzzeitig durch die Schaltung und der Kondensator entlädt sich dabei. Der Strom I als auch die Spannung U-C nehmen mit einer e-Funktion ab und gehen gegen 0.
		Warum entläd sich der Kondensator? Der Strom der durch die Schaltung fießt wird von der Spannung U-C also der Ladungsdifferenz der Platten getrieben.
		Die Frage die sich hierbei stellt ist, wie lange dauert dieser entlade oder eben auflade Vorgang. Hierfür wird Tau verwendet.
		Ein Tau entspricht der Zeit, bis die Kondensatorspannung um 63 Prozent ihres Ursprungswert beim entladen abfällt. Beim Aufladen wird die Zeit, bis die Spannung auf 63 Prozent ihres Endwert angestiegen ist, verwendet.
		Diese Zeit kann auch aus der Kapazität und dem Serienwiderstand mit Tau ist gleich R mal C berechnet werden.
	}

\end{frame}


%%%%% Laden eines Kondensators
\s{
	\begin{figure}[H]
		\centering
		\begin{subfigure}[h]{0.9\textwidth}
			\centering
			\input{Tikz/src/03_Kapazitaet_Kondensator/Laden_Kondensator_ESB_Schalter}
			\caption{Darstellung der Reihenschaltung bestehend aus einer Kapazität \(C\), einem Widerstand \(R\) und einem Schalter.}
			\label{fig:Schaltbild_CR}
		\end{subfigure}

		%%%%%%%%%%%%%%%%%%%%%%%%%%%%%%%%%%%%%%%%%%%%%%%%% Ende - Schaltbild mit Schalter	 %%%%%%%%%%%%%%%%%%%%%%%%%%%%%%%%%%%%%%%%%%%%%%
		%%%%%%%%%%%%%%%%%%%%%%%%%%%%%%%%%%%%%%%%%%%%%%%%% Anfang - Diagramm mit Schalter	 %%%%%%%%%%%%%%%%%%%%%%%%%%%%%%%%%%%%%%%%%%%%%%
		\begin{subfigure}[b]{0.9\textwidth}
			\centering
			\input{Tikz/src/03_Kapazitaet_Kondensator/Spannungs_Stromverlauf_Kapazitaet_Skript}
			\caption{Spannungs- und Strom verlauf über der Kapazität bei einer Gleichspannung die zum Zeitpunkt $t = 0$ eingeschaltet wird und zum Zeitpunkt $t_\mathrm{0}$ abgeschaltet wird.}
			\label{fig:Spannungsverlauf_CR_1}
		\end{subfigure}
		\caption{\textbf{Der Spannungs- und Stromverlauf über der Kapazität.} Das Diagramm zeigt die Reaktion der Kapazität auf die plötzliche Änderung der Spannung und
			veranschaulicht das typische Lade verhalten.}
		\alttext{Zweiteilige Abbildung zum Laden eines Kondensators in einer Reihenschaltung. Oben ist eine Schaltung mit Gleichspannungsquelle Uq, einem Schalter, einem Kondensator C und einem Widerstand R in Reihe dargestellt. Der Schalter wird bei t0 gleich 5 Tau geöffnet. Pfeile markieren die Spannungen uC von t am Kondensator und uR von t am Widerstand sowie den Strom i von t durch den Zweig. Unten ist ein Zeitdiagramm gezeigt: die Quellspannung Uq springt bei t gleich 0 auf 100 Prozent und bleibt bis t0 konstant, danach fällt sie auf null. Die Kondensatorspannung uC von t steigt exponentiell von 0 auf nahezu 100 Prozent an. Der Strom i von t startet bei 100 Prozent und fällt exponentiell gegen 0. Bei t gleich Tau sind 63 Prozent für uC und 37 Prozent für i markiert.}
		\label{fig:Gesamtdarstellung_1}
	\end{figure}
}

%%%%%%%%%%%%%%%%%%%%%%%%%%%%%%%%%%%%%%%%%%%%%%%%% Ende - Diagramm mit Schalter	 %%%%%%%%%%%%%%%%%%%%%%%%%%%%%%%%%%%%%%%%%%%%%%
\s{
	Im Diagramm \ref{fig:Spannungsverlauf_CR_1} ist zu sehen, dass sich trotz angelegter Gleichspannung $U_\mathrm{q}$ die Spannung $u_\mathrm{c}(t)$
	zeitgleich nicht mitändert, sondern eine Trägheit aufweisen. Dieses zeitlich veränderliche Verhalten wird im Schaltbild \ref{fig:Schaltbild_CR}
	durch die zeitabhängigen Größen $u(t)$ und $i(t)$ ausgedrückt.

	In diesem Beispiel ist der Schalter zu Beginn geschlossen und wird bei $t_\mathrm{0} = 5\tau$ geöffnet.
	Der Kondensator ist anfangs nicht geladen. Im Diagramm ist zu erkennen, dass sich trotzt angelegter Spannung $U_\mathrm{q}$
	die Spannung über der Kapazität $u_\mathrm{C}(t)$ nicht sprunghaft ändert.

	Dieses Einschaltverhalten resultiert aus der Aufladung der Kapazität. Die Ladungsträger müssen sich dafür auf den Kondensatorplatten
	ansammeln. Dies geschieht so lange, bis die Spannung zwischen den Platten der angelegten Spannung der Spannungsquelle entspricht.
	Da die zur Aufladung benötigten Ladungsträger in die Kondensatorplatten wandern müssen, fließt während dieses Vorgangs ein entsprechender Strom $i_\mathrm{C}(t)$.

	Im Fall einer ideal leitenden Schaltung ohne ohm'schen Widerstand, würden sich die Ladungsträger unendlich schnell auf den Kondensatorplatten ansammeln,
	was einem unendlich hohen Strom entsprechen würde. Im realen Fall hat der Kondensator bzw. die Hinleitung zum Kondensator, einen ohm'schen Widerstand,
	welcher die Ladungsträger auf dem Weg zu den Kondensatorplatten ausbremst, also die elektrische Stromstärke reduziert.
	Daraus folgt der maximal mögliche Strom $I_\mathrm{max}$.

	\begin{equation}
		I_\mathrm{max} = \frac{U_\mathrm{q}}{R}
	\end{equation}

	% Da es dabei dennoch zu relativ hohen Stromstärken kommen kann, werden in Schaltungen mit Kondensatoren mit hohen Kapazitäten sogenannten Vorladewiderstand verbaut.
	% Deren Aufgabe ist es, den maximalen Strom zu begrenzen.

	Sobald der Schalter geöffnet wird, bleibt die Spannung über der Kapazität erhalten.
	Dies bedeutet, dass die zugeführte Energie in der Kapazität verbleibt und sich im stationären Fall nicht mehr ändert.
	Dementsprechend fließt bei einer aufgeladenen Kapazität auch kein Strom mehr, was in einem eingeschwungenen Gleichstromnetzwerk einem Leerlauf gleichzusetzen ist.

	Mit Hilfe der beiden Formeln \ref{eq:ict} und \ref{eq:uci} ist das Verhalten von Strom und Spannung an einer Kapazität vollständig beschrieben,
	was nun diverse Berechnungen ermöglicht. Ein weiteres Beispiel ist das Schaltverhalten einer Kapazität bei anlegen einer Rechteckspannung.
	Der Unterschied besteht darin, dass nicht wie in Beispiel \ref{fig:Schaltbild_CR}, eine konstante Spannungsquelle über einen Schalter geschaltet wird,
	sondern die Spannungsquelle selbst die Eingangsspannung bestimmt.
}
%%%%%%%%%%%%%%%%%%%%%%%%%%%%%%%%%%%%%%%%%%%%%%%%% Anfang - Schaltbild mit Rechteck	 %%%%%%%%%%%%%%%%%%%%%%%%%%%%%%%%%%%%%%%%%%%%%%
\iffalse
	\begin{frame}[t]
		\ftx{Schaltverhalten eines Kondensators: Entladen}

		% ; TODO: Hillgärtner: 	Da habe ich so mein Problem mit, weil diese Rechteckspannung in dieser Spannungsquelle würde ja heißen, dass die Spannung, wenn sie auf 0 Volt in der Spannungsquelle fällt, quasi ein Kurzschluss ist. Es führt immer zu Diskussionen. Ich würde es anders machen. Ich würde die Spannungsquelle durch einen Kurzschluss ersetzen, einen Schalter einbauen, der zum Zeitpunkt T gleich 0 geschlossen wird. Und sagen, dass die Anfangsbedingung ist, dass die Kapazität C auf eine Spannung UQ aufgeladen ist. Und dann kann man die Entladekurve aufzeichnen und den Stromfluss aufzeichnen. Das fände ich sinnvoll. Alles andere finde ich immer sehr schwierig zu erklären mit diesen Rechteckspannungsquellen, weil in der Realität ist es nicht unbedingt immer so.
		%						Demzufolge müssten die Folien 98, 99, 100, 101, 102, 103 angepasst werden. Ich weiß nicht, ob das mit der Rechteckspannung irgendwann nochmal in weiteren Modulen kommt, aber wir sollen ja nur die Kapazität prinzipiell erläutern. Und ich finde, da reichen die Kurven völlig aus.

		%wurde neu gestatltet durch JW

		\b{
			%%%%% Laden- entladen eines Kondensators
			\begin{figure}[H]
				\centering
				\begin{subfigure}[h]{1\textwidth}
					\centering
					\input{Tikz/src/03_Kapazitaet_Kondensator/Kapazitaet_Reihenschaltung_Entladen}
				\end{subfigure}
				%%%%%%%%%%%%%%%%%%%%%%%%%%%%%%%%%%%%%%%%%%%%%%%%% Ende - Schaltbild mit Rechteck	 %%%%%%%%%%%%%%%%%%%%%%%%%%%%%%%%%%%%%%%%%%%%%%
				%%%%%%%%%%%%%%%%%%%%%%%%%%%%%%%%%%%%%%%%%%%%%%%%% Anfang - Diagramm mit Rechteck	 %%%%%%%%%%%%%%%%%%%%%%%%%%%%%%%%%%%%%%%%%%%%%%
				\only<2->{
					\begin{subfigure}[b]{1\textwidth}
						\centering
						\input{Tikz/src/03_Kapazitaet_Kondensator/Spannungs_Stromverlauf_Kapazitaet_Entladen_Folien}
					\end{subfigure}}
			\end{figure}

		}
	\end{frame}
\fi
\s{
	%%%%% Laden- entladen eines Kondensators
	\begin{figure}[H]
		\centering
		\begin{subfigure}[h]{0.85\textwidth}
			\centering
			\input{Tikz/src/03_Kapazitaet_Kondensator/Kapazitaet_Reihenschaltung_Entladen}
			\caption{Darstellung der Reihenschaltung bestehend aus einer Kapazität \(C\) und einem Widerstand \(R\) an einer rechteckförmigen Spannungsquelle.}
			\label{fig:Schaltbild_CRR}
		\end{subfigure}
		%%%%%%%%%%%%%%%%%%%%%%%%%%%%%%%%%%%%%%%%%%%%%%%%% Ende - Schaltbild mit Rechteck	 %%%%%%%%%%%%%%%%%%%%%%%%%%%%%%%%%%%%%%%%%%%%%%
		%%%%%%%%%%%%%%%%%%%%%%%%%%%%%%%%%%%%%%%%%%%%%%%%% Anfang - Diagramm mit Rechteck	 %%%%%%%%%%%%%%%%%%%%%%%%%%%%%%%%%%%%%%%%%%%%%%
		\begin{subfigure}[b]{0.9\textwidth}
			\centering
			\input{Tikz/src/03_Kapazitaet_Kondensator/Spannungs_Stromverlauf_Kapazitaet_Entladen_Skript}
			\caption{Spannungs- und Strom verlauf über der Kapazität bei einer rechteckförmigen Spannung.}
			\label{fig:Spannungsverlauf_2}
		\end{subfigure}
		\caption{\textbf{Spannungs- und Stromverlauf über der Kapazität.} Das Diagramm zeigt die Reaktion der Kapazität auf eine plötzliche Änderung
			der Spannung und veranschaulicht das typische Lade- und Entladeverhalten des Kondensators.}
			\alttext{Zweiteilige Abbildung zum Lade- und Entladeverhalten eines RC-Glieds bei rechteckförmiger Anregung. Oben ist eine Reihenschaltung aus Kondensator C und Widerstand R an einer rechteckförmigen Spannungsquelle u von t dargestellt; Pfeile markieren die Kondensatorspannung uC von t, die Widerstandsspannung uR von t und den Strom i von t. Unten zeigt ein Zeitdiagramm drei Verläufe: die Eingangsspannung u von t springt auf 100 Prozent und fällt bei 5 Tau auf 0. Die Kondensatorspannung uC von t steigt bis 5 Tau exponentiell gegen 100 Prozent und fällt danach exponentiell wieder gegen 0. Der Strom i von t startet positiv und fällt gegen 0; beim Abschalten wird er negativ und nähert sich wieder 0. Markierungen zeigen bei Tau etwa 63 Prozent für uC und 37 Prozent für i; nach dem Abschalten ist bei 6 Tau ein Punkt auf uC sowie ein negativer Stromwert markiert.}
		\label{fig:Gesamtdarstellung_2}
	\end{figure}
}% Nur Skript

%%%%%%%%%%%%%%%%%%%%%%%%%%%%%%%%%%%%%%%%%%%%%%%%% Ende - Diagramm mit Rechteck	 %%%%%%%%%%%%%%%%%%%%%%%%%%%%%%%%%%%%%%%%%%%%%%
\s{
	Im Gegensatz zum Schalter, versucht die Spannungsquelle im Beispiel \ref{fig:Schaltbild_CRR}
	ab dem Zeitpunkt $t_\mathrm{0} = 5\tau$ die Systemspannung auf $0$ zu ziehen. Dies führt dazu,
	dass die Kapazität entladen wird. Dies führt entsprechend zu einem entgegengesetzten Stromfluss $i_\mathrm{C}(t)$,
	da sich die Ladungsträger von den Kondensatorplatten wegbewegen, um wieder einen neutralen Zustand zu erreichen.
	%Neben der Verwendung zur Energiespeicherung, werden Kondensatoren auch für die Filterung von Frequenzen, zur Zeitmessung oder für die Signalverarbeitung genutzt.
}
%%%%%%%%%%%%%%%%%%%%%%%%%%%%%%%%%%%%%%%%%%%%%%%%% Ende - Schaltverhalten Kondensator	 %%%%%%%%%%%%%%%%%%%%%%%%%%%%%%%%%%%%%%%%%%%%%%
%%%%%%%%%%%%%%%%%%%%%%%%%%%%%%%% Anfang - Beispiel Rechnung Kondensator %%%%%%%%%%%%%%%%%%%%%%%%%%%%%%%

% \begin{frame}
% 	\ftx{Beispielrechnung Kapazität}

% 	\begin{bsp}{Berechnung der Kapazität}{}
% 		\begin{itemize}

% 			\item Ladung auf Platten $Q = \sigma \cdot A = \varepsilon \cdot E \cdot A$
% 			\item Kapazität mit $C = \frac{\varepsilon \cdot A}{d}$ berechnen
% 			\item Gespeicherte Energie $W = \frac{1}{2} \cdot C \cdot U^2 $
% 			\item Kapazität zweier Kondensatoren mit verschiedenen $\varepsilon_{\mathrm{r}}$ über $\vec{D}$ berechnen

% 		\end{itemize}
% 	\end{bsp}

% \end{frame}
%%%%%%%%%%%%%%%%%%%%%%%%%%%%%%%% Ende - Beispiel Rechnung Kondensator %%%%%%%%%%%%%%%%%%%%%%%%%%%%%%%
